\documentclass{article}
\usepackage{amsmath, amssymb, amsthm, geometry, hyperref}
\geometry{a4paper, margin=1in}

\title{ACM India Summer School on Symmetric Key Cryptography \\ Day 6: Analysis and Implementation of the Salsa20 Stream Cipher}
\author{Sourav Sharma}
\date{June 13, 2026}

\begin{document}

\maketitle

\section{Introduction}
Salsa20 is a high-speed stream cipher designed by Daniel J. Bernstein, based on the Add-Rotate-XOR (ARX) operations. The cipher operates on a state of sixteen 32-bit words (512 bits in total), arranged in a $4 \times 4$ grid. In this lab, I implemented the core Salsa20 transformations, encryption/decryption mechanisms, and studied its differential propagation properties.

\section{Salsa20 Core Functions}

\subsection{Quarter Round (QR)}
The Quarter Round function is the fundamental ARX building block of Salsa20. It operates on four 32-bit words $(a, b, c, d)$ and updates them in place:
\begin{align*}
    b &\gets b \oplus \text{rotl}(a + d, 7) \\
    c &\gets c \oplus \text{rotl}(b + a, 9) \\
    d &\gets d \oplus \text{rotl}(c + b, 13) \\
    a &\gets a \oplus \text{rotl}(d + c, 18)
\end{align*}
where $\text{rotl}(x, n)$ represents a left circular shift of $x$ by $n$ bits.

\subsection{Row, Column, and Double Rounds}
Salsa20 alternates rounds of column and row transformations. A Column Round applies the QR function to the four columns of the state grid:
\[ \text{qr}(0, 4, 8, 12), \quad \text{qr}(5, 9, 13, 1), \quad \text{qr}(10, 14, 2, 6), \quad \text{qr}(15, 3, 7, 11) \]
A Row Round applies the QR function to the rows (with diagonal starting points):
\[ \text{qr}(0, 1, 2, 3), \quad \text{qr}(5, 6, 7, 4), \quad \text{qr}(10, 11, 8, 9), \quad \text{qr}(15, 12, 13, 14) \]
A Double Round is defined as one Column Round followed by one Row Round. Salsa20/20 executes 10 double rounds (20 rounds total).

\section{Salsa20 Encryption and Decryption}
The 512-bit state matrix is initialized with four constant words, an 8-word (256-bit) key, a 2-word (64-bit) nonce, and a 2-word (64-bit) block counter:
\[
\begin{pmatrix}
    \text{const}_0 & \text{key}_0 & \text{key}_1 & \text{key}_2 \\
    \text{key}_3 & \text{const}_1 & \text{nonce}_0 & \text{nonce}_1 \\
    \text{ctr}_0 & \text{ctr}_1 & \text{const}_2 & \text{key}_4 \\
    \text{key}_5 & \text{key}_6 & \text{key}_7 & \text{const}_3
\end{pmatrix}
\]
The keystream block is generated by running 10 double rounds on the state and then adding the original input state word-by-word to the final round output (feed-forward step to ensure invertibility is impossible without the key).
Encryption and decryption are symmetric; both are performed by XORing the plaintext/ciphertext with the generated keystream. I verified my implementation (\href{file:///home/scorzion/Tech/Crypto/SKC_IITH/Day6/solutions/salsa_enc_dec.cpp}{\texttt{salsa\_enc\_dec.cpp}}) on a test string, achieving full decryption recovery.

\section{Differential Cryptanalysis}

\subsection{Hamming Weight Propagation after One Round}
We studied the propagation of a single-bit input difference at $\Delta X = X_7[31]$ (MSB of the 8th word) through 1 column round (which is the first round of a double round).
Because columns 0, 1, and 2 are inactive (having zero difference), only Column 3 (indices 15, 3, 7, 11) is active:
\begin{enumerate}
    \item $b \gets 3, a \gets 15, d \gets 11$: Since $X_{15}$ and $X_{11}$ have zero difference, $X_3$ receives zero difference.
    \item $c \gets 7, b \gets 3, a \gets 15$: Since $X_3$ and $X_{15}$ have zero difference, the addition/rotation yields zero difference. Thus, $X_7$ retains its original input difference: $\Delta X_7 = 2^{31}$ (Hamming weight = 1).
    \item $d \gets 11, c \gets 7, b \gets 3$: Here we compute $\text{rotl}(X_7 + X_3, 13)$. Since the input difference is at the MSB of $X_7$ ($2^{31}$), addition modulo $2^{32}$ preserves the MSB difference exactly (no carry is possible beyond the 31st bit). Thus, the difference entering the rotation is exactly $2^{31}$. Rotating left by 13 bits gives a difference at bit 12: $\Delta X_{11} = 2^{12}$ (Hamming weight = 1).
    \item $a \gets 15, d \gets 11, c \gets 7$: We compute $\text{rotl}(X_{11} + X_7, 18)$. The input difference in $X_{11}$ is at bit 12, which propagates due to carries in addition, resulting in an average carry chain length of 2 (affecting bits 12, 13, etc.). The difference in $X_7$ is at bit 31. Therefore, the sum difference has an average Hamming weight of $\approx 3$ active bits. Rotating this by 18 bits preserves the Hamming weight. Thus, $\Delta X_{15}$ receives an average difference weight of $\approx 3$.
\end{enumerate}
Adding these active differences, the total average Hamming weight of the state difference after 1 round is:
\[ \text{HW}(\Delta Y) = \text{HW}(\Delta X_7) + \text{HW}(\Delta X_{11}) + \text{HW}(\Delta X_{15}) \approx 1 + 1 + 3 = 5 \]
This theoretical prediction is verified empirically by my simulation in \href{file:///home/scorzion/Tech/Crypto/SKC_IITH/Day6/solutions/Salsa20.cpp}{\texttt{Salsa20.cpp}}, which reports an average Hamming weight of exactly \textbf{4.985} over $100,000$ random samples.

\subsection{Differential Probability of Salsa20/4}
We estimated the differential probability for the transition $\Delta X = X_7[31] \to \Delta Y = Y_1[14]$ after 4 rounds (Salsa20/4). Over $2^{22}$ simulated random pairs, the probability of no difference ($\Delta Y_1[14] = 0$) was measured to be approximately \textbf{0.5412}. This shows a significant bias away from the uniform distribution value of $0.50$, presenting a potential distinguisher.

\end{document}
